% ─────────────────────────────────────────────────────────────────────────────
%  中文段落格式：首行缩进 + 标题后首段也缩进
% ─────────────────────────────────────────────────────────────────────────────
%  为什么需要这个文件
%    Eisvogel 默认走 KOMA-Script 的 `parskip=half`：**首行不缩进、段间空一行**。
%    那是西文排版的做法。中文的惯例是「段间不空行、每段首行缩进」。
%    构建脚本传了 `-V indent=true` 让 Eisvogel **跳过**那段设置，
%    再由这里把 \parindent 显式定成 1em（10pt 字号下正好是一个汉字宽）。
%
%  为什么还要 indentfirst
%    LaTeX 的 `\@startsection` 默认 `\@afterindentfalse` —— 紧跟标题的
%    **第一段不缩进**。中文排版要求它也缩进，所以加载 indentfirst
%    （它做的就是把 \@afterindentfalse 换成 \@afterindenttrue）。
%    《中文信息学报》等模板同样加载了这个宏包。
%
%  ⚠ `@` 在 pandoc 的 header-includes 里不是字母（catcode 12），
%    所以不在这里直接写 \@afterindenttrue，交给 indentfirst 宏包处理。
% ─────────────────────────────────────────────────────────────────────────────

\usepackage{indentfirst}

% 用 AtBeginDocument 兜底：确保最终生效的是 1em，
% 不受宏包加载顺序或 KOMA 默认值的影响。
\AtBeginDocument{%
  \setlength{\parindent}{1em}%
  \setlength{\parskip}{0pt plus 0.5pt}%
}
